\subsection{Точки разрыва функции и их классификация, примеры.}

Пусть $f(x)$ определена в проколотой окрестности $\mathring{U}_a$ (в самой точке $a$ может быть не определена). \\
Точка $a$ называется \textbf{точкой разрыва}, если $f(x)$ не является непрерывной в этой точке (не определена, предел не существует или предел не равен значению).

\begin{enumerate}
	\item \textbf{Разрывы первого рода} (оба односторонних предела существуют и конечны):
	\begin{enumerate}
		\item \textbf{Устранимый разрыв:} $\lim_{x\to a} f(x)$ существует, но либо $f(a)$ не существует, либо $\lim_{x\to a} f(x) \neq f(a)$.
		\item \textbf{Разрыв первого рода (скачок):} $\lim_{x\to a+0} f(x) \neq \lim_{x\to a-0} f(x)$.
	\end{enumerate}

	\item \textbf{Разрывы второго рода:}
		Хотя бы один из односторонних пределов не существует или равен бесконечности.
\end{enumerate}

\textbf{Примеры:}
\begin{enumerate}
	\item $f(x) = \operatorname{sgn} x$. В точке $x=0$: пределы $\pm 1$. \textbf{Скачок (I род).}
	\item $f(x) = \operatorname{sgn}^2 x$ (равна 1 везде, кроме 0; в нуле 0). В точке $x=0$: предел 1, значение 0. \textbf{Устранимый (I род).}
	\item $f(x) = \sin \frac{1}{x}$. В точке $x=0$ предел не существует (осцилляция). \textbf{II род.}
	\item $f(x) = \frac{1}{x}$. В точке $x=0$ пределы бесконечны. \textbf{II род (бесконечный).}
\end{enumerate}

